\chapter{2014:Eric Betzig and Stefan W. Hell and William E. Moerner}

\label{chap:chemistry2014}

The Nobel Prize in Chemistry for the year 2014 was awarded jointly to Eric Betzig and Stefan W. Hell and William E. Moerner.

\textbf{Motivation:}

for the development of super-resolved fluorescence microscopy

\section{Eric Betzig}

\subsection*{Early Life and Education}

Eric Betzig was born in 1960 in Ann Arbor, Michigan, USA.

\subsection*{Career and Major Research}

Betzig earned a bachelor's degree from the California Institute of Technology in 1983 and a doctorate in applied physics from Cornell University in 1988. He worked at Bell Laboratories and later at the Janelia Research Campus of the Howard Hughes Medical Institute, where he developed photoactivated localization microscopy (PALM) together with Harald Hess.

\subsection*{Nobel Prize-winning Work}

The work for which Eric Betzig was awarded the Nobel Prize in Chemistry in 2014 was described by the Nobel Committee as: "for the development of super-resolved fluorescence microscopy".

PALM switches on fluorescent molecules one at a time and records their positions, building up a high-resolution image of a cell from many individual localizations, far below the diffraction limit of light.

\subsection*{Significance and Impact}

Betzig's method allows cellular structures to be imaged at nanometer resolution, revealing details that optical microscopy could not previously resolve.

\subsection*{Later Career and Honors}

Betzig continues his research on high-resolution imaging at the Janelia Research Campus.

\subsection*{Legacy}

PALM and related single-molecule methods have made nanoscale imaging of living cells possible and are used across cell biology.

\subsection*{Selected Publications}

"Imaging Intracellular Fluorescent Proteins at Nanometer Resolution" (Science, 2006)

\clearpage

\section{Stefan W. Hell}

\subsection*{Early Life and Education}

Stefan W. Hell was born in 1962 in Arad, Romania.

\subsection*{Career and Major Research}

Hell studied physics and received his doctorate from the University of Heidelberg in 1990. He became a director at the Max Planck Institute for Biophysical Chemistry in Göttingen, where he developed stimulated emission depletion (STED) microscopy.

\subsection*{Nobel Prize-winning Work}

The work for which Stefan W. Hell was awarded the Nobel Prize in Chemistry in 2014 was described by the Nobel Committee as: "for the development of super-resolved fluorescence microscopy".

STED microscopy uses a second laser beam to switch off fluorescence everywhere except at the center of the spot, shrinking the effective focus far below the diffraction limit.

\subsection*{Significance and Impact}

STED gave scientists a way to image structures inside cells with resolution beyond the classical limit of light microscopy, and it works on living samples.

\subsection*{Later Career and Honors}

Hell continues his research on super-resolution microscopy at the Max Planck Institute for Biophysical Chemistry in Göttingen.

\subsection*{Legacy}

STED microscopy is now a commercialized and widely used technique for nanoscale imaging in the life sciences.

\subsection*{Selected Publications}

"Fluorescence Microscopy with Diffraction Resolution Barrier Broken by Stimulated Emission" (Proceedings of the National Academy of Sciences, 2000)

\clearpage

\section{William E. Moerner}

\subsection*{Early Life and Education}

William E. Moerner was born in 1953 in Pleasanton, California, USA.

\subsection*{Career and Major Research}

Moerner earned a bachelor's degree from Washington University in St. Louis in 1975 and a doctorate from Cornell University in 1982. He worked at IBM Almaden Research Center and later became a professor at Stanford University, where he studied the optics of single molecules.

\subsection*{Nobel Prize-winning Work}

The work for which William E. Moerner was awarded the Nobel Prize in Chemistry in 2014 was described by the Nobel Committee as: "for the development of super-resolved fluorescence microscopy".

Moerner achieved the first optical detection and spectroscopy of a single molecule in a solid, demonstrating that individual molecules could be observed and studied.

\subsection*{Significance and Impact}

His single-molecule measurements showed that individual molecules differ in their behavior, and they provided the physical basis for the localization methods used in super-resolution microscopy.

\subsection*{Later Career and Honors}

Moerner continues his research on single-molecule imaging and spectroscopy at Stanford University.

\subsection*{Legacy}

Moerner's single-molecule techniques opened the field of single-molecule chemistry and are a foundation of modern super-resolution imaging.

\subsection*{Selected Publications}

"Optical Detection and Spectroscopy of Single Molecules in a Solid" (Physical Review Letters, 1989)

\clearpage

\section*{Chapter Summary}

The 2014 prize recognized three methods that break the diffraction barrier of light microscopy: STED by Hell, PALM by Betzig, and the single-molecule detection by Moerner on which localization methods rely. Together they allow biological structures to be imaged at the nanometer scale.
